 vytvářejí specifické, vysoce fragmentované segmenty trhu práce, do nichž dosavadní standardní mechanismy kolektivního vyjednávání v~\acloc{ČR} zasahují jen zcela okrajově. Tento vývoj s~sebou nese strukturální riziko vzniku dvourychlostního trhu práce: zatímco úzké jádro vysoce kvalifikovaných zaměstnanců v~silných podnicích disponuje mzdovou ochranou, rozsáhlé periferní segmenty v~nízkokvalifikovaných službách, logistických skladech, stavebnictví a~ve stavebních subdodavatelských řetězcích jsou ponechány bez jakýchkoliv kolektivních standardů, což vytváří tlak na celkové stlačování tuzemských mezd.\par

Masivní ukrajinská migrační vlna po roce 2022 vykazuje v~tomto kontextu rozporuplný a~asymetrický efekt. Na jedné straně sice významným způsobem pomohla stabilizovat demografickou a~pracovní bilanci českého trhu práce v~době převisu volných pracovních míst, avšak na straně druhé ukazuje na selhání integračních mechanismů. Velká část vysoce kvalifikovaných uprchlíků končí na pozicích hluboko pod úrovní svého vzdělání, v~šedé ekonomice, nebo v~prekérním agenturním zaměstnávání se minimální mzdovou ochranou. Pakliže integrační politika a~státní správa efektivně nezapojují sociální partnery do tvorby mzdových a~pracovních standardů pro nově příchozí, hrozí vznik odděleného, legislativně i~mzdově podřadného segmentu práce, což eliminuje tlak na technologickou modernizaci a~přirozené zvyšování mezd u~domácích nízkopříjmových profesí.~\cite{PolentsDvorakova_UkraineRefugees2024}\par

Digitální platformová práce posouvá toto napětí do virtuálního prostoru a~fakticky dekonstruuje samotný tradiční koncept zaměstnaneckého poměru. Návrhy české právní úpravy a~nová evropská směrnice z~roku 2024 o~platformové práci přinášejí přelomové zpřísnění v~podobě vyvratitelné domněnky existence pracovního poměru. Výzvou zůstává otázka reprezentativnosti: pokud tito zranitelní a~digitálně řízení pracovníci zůstanou vyčleněni z~působnosti odvětvových standardů, bude se dosah sociálního dialogu v~\acloc{ČR} zužovat na stále menší a~izolovanější skupinu tradičních zaměstnanců, což povede k~další erozi celospolečenské vyjednávací síly. Společné působení migračních tlaků a~platformové práce tak uzavírá přetrvávající model levné práce, nucené participace a~nízkých standardů na českém trhu práce.\par
